\documentclass[../../main.tex]{subfiles}

\begin{document}

\section{笛卡儿积}\label{section:ch01-笛卡儿积}

我们已经定义了两个集合的笛卡儿积 $A \times B$。现在介绍更为一般的笛卡儿积。

\begin{definition}[加标集族]\label{definition:ch01-加标集族}
设 $\mathcal{A}$ 是一个非空集族。$\mathcal{A}$ 的指标函数（indexing function）是从某一个集合 $J$ 到 $\mathcal{A}$ 的一个满射 $f$，其中 $J$ 称为指标集（index set）。族 $\mathcal{A}$ 连同指标函数 $f$ 一起称为一个集的加标族（indexed family of sets）或加标集族（indexed family of sets）。给定 $\alpha \in J$，集合 $f(\alpha)$ 记成符号 $A_{\alpha}$。这个加标集族本身则记作
\begin{align*}
\{A_{\alpha}\}_{\alpha \in J},
\end{align*}
读作“$\alpha$ 取遍 $J$ 时，所有 $A_{\alpha}$ 的族”。当指标集自明时，则简单地记为 $\{A_{\alpha}\}$。
\end{definition}

注意，虽然要求指标函数为满射，但它不一定要是单射，甚至对于所有的 $\alpha \neq \beta$，$A_{\alpha}$ 与 $A_{\beta}$ 是 $\mathcal{A}$ 中的同一个集合也是完全可以的。

指标函数的用处之一便是可以给集合的任意并与交一个新的记号。设 $f: J \longrightarrow \mathcal{A}$ 是 $\mathcal{A}$ 的一个指标函数，并且用 $A_{\alpha}$ 表示 $f(\alpha)$，则定义
\begin{align}
\bigcup_{\alpha \in J} A_{\alpha} = \{x \mid \text{对于至少一个 } \alpha \in J, x \in A_{\alpha}\},\label{align:ch01-加标任意并}
\end{align}
以及
\begin{align}
\bigcap_{\alpha \in J} A_{\alpha} = \{x \mid \text{对于任意 } \alpha \in J, x \in A_{\alpha}\}.\label{align:ch01-加标任意交}
\end{align}
这就是上面定义的概念的新的简单记号，（由于指标函数是满射）易见，第一个式子是 $\mathcal{A}$ 中所有元素的并，第二个式子是 $\mathcal{A}$ 中所有元素的交。

从 $1$ 到 $n$ 的正整数的集合 $\{1, \cdots, n\}$ 和正整数集 $\mathbb{Z}_{+}$ 是两个极为常用的指标集。对于这两个指标集，我们引入几个常用记号。以 $\{1, \cdots, n\}$ 为指标集的加标集族记为 $\{A_{1}, \cdots, A_{n}\}$，其成员的并和交分别记作：
\begin{align*}
A_{1} \cup \cdots \cup A_{n} \quad \text{和} \quad A_{1} \cap \cdots \cap A_{n}.
\end{align*}
对于指标集为 $\mathbb{Z}_{+}$ 情形，加标集族记为 $\{A_{1}, A_{2}, \cdots\}$，其成员的并和交分别记作：
\begin{align*}
A_{1} \cup A_{2} \cup \cdots \quad \text{和} \quad A_{1} \cap A_{2} \cap \cdots.
\end{align*}

\begin{definition}[$m$-串]\label{definition:ch01-m串}
设 $m$ 是一个正整数。对于给定的集合 $X$，$X$ 中元素的一个 $m$-串（$m$-tuple）定义为函数
\begin{align*}
\boldsymbol{x}: \{1, \cdots, m\} \longrightarrow X.
\end{align*}
若 $\boldsymbol{x}$ 是一个 $m$-串，则往往把 $\boldsymbol{x}$ 在 $i$ 处的值写成符号 $x_{i}$ 而不记成 $\boldsymbol{x}(i)$，并称之为 $\boldsymbol{x}$ 的第 $i$ 个坐标（coordinate），函数 $\boldsymbol{x}$ 本身用符号
\begin{align*}
(x_{1}, \cdots, x_{m})
\end{align*}
表示。
\end{definition}

现在设 $\{A_{1}, \cdots, A_{m}\}$ 是一个以 $\{1, \cdots, m\}$ 为指标集的加标集族。令 $X = A_{1} \cup \cdots \cup A_{m}$。这个加标族的笛卡儿积（Cartesian product），记作
\begin{align}
\prod_{i=1}^{m} A_{i} \quad \text{或者} \quad A_{1} \times \cdots \times A_{m},\label{align:ch01-m重笛卡儿积}
\end{align}
定义为 $X$ 中元素的所有 $m$-串 $(x_{1}, \cdots, x_{m})$ 的集合，使得对于每一个 $i$ 有 $x_{i} \in A_{i}$。

\begin{example}\label{example:ch01-两种笛卡儿积}
符号 $A \times B$ 现在有了两种定义。当然有一种定义是早先给出的，它用 $A \times B$ 表示所有有序偶对 $(a, b)$ 的集合，其中 $a \in A$，$b \in B$。第二种定义是刚才给出的，将 $A \times B$ 定义为所有函数 $\boldsymbol{x}: \{1, 2\} \longrightarrow A \cup B$ 的集合，其中 $x(1) \in A$，$x(2) \in B$。在这两个集合之间有一个明显的一一对应：有序偶对 $(a, b)$ 对应于由 $x(1) = a$，$x(2) = b$ 定义的函数 $\boldsymbol{x}$。因为用“串记号”这种函数 $\boldsymbol{x}$ 通常记成符号 $(a, b)$，这本身就给出了对应。因此，两个集合的笛卡儿积的这种一般定义，从本质上说就是早先的那种定义。
\end{example}

\begin{example}\label{example:ch01-结合性}
笛卡儿积 $A \times B \times C$ 与笛卡儿积 $A \times (B \times C)$，$(A \times B) \times C$ 之间有多少差别呢？差别很小。因为它们之间存在着明显的一一对应：
\begin{align*}
(a, b, c) \longleftrightarrow (a, (b, c)) \longleftrightarrow ((a, b), c).
\end{align*}
\end{example}

\begin{definition}[$\omega$-串]\label{definition:ch01-omega串}
给定一个集合 $X$，定义 $X$ 中元素的 $\omega$-串（$\omega$-tuple）为函数
\begin{align*}
\boldsymbol{x}: \mathbb{Z}_{+} \longrightarrow X;
\end{align*}
这种函数也称为 $X$ 中元素的一个序列（sequence）或一个无穷序列（infinite sequence）。若 $\boldsymbol{x}$ 是一个 $\omega$-串，则往往把 $\boldsymbol{x}$ 在 $i$ 处的值写成 $x_{i}$ 而不写成 $\boldsymbol{x}(i)$，并称之为 $\boldsymbol{x}$ 的第 $i$ 个坐标（coordinate），$\boldsymbol{x}$ 本身用符号
\begin{align*}
(x_{1}, x_{2}, \cdots) \quad \text{或者} \quad (x_{n})_{n \in \mathbb{Z}_{+}}
\end{align*}
表示。设 $\{A_{1}, A_{2}, \cdots\}$ 是以正整数集作为指标集的一个加标集族，$X$ 为这个集族中所有集合之并。这个加标集族的笛卡儿积（Cartesian product）记成
\begin{align}
\prod_{i \in \mathbb{Z}_{+}} A_{i} \quad \text{或者} \quad A_{1} \times A_{2} \times \cdots,\label{align:ch01-无穷笛卡儿积}
\end{align}
定义为 $X$ 中元素的所有 $\omega$-串 $(x_{1}, x_{2}, \cdots)$ 的集合，使得对于每一个 $i$，$x_{i} \in A_{i}$。
\end{definition}

上述定义中，并不要求 $A_{i}$ 两两不同。事实上，它们可以取同一个集合 $X$。对于这种情形，笛卡儿积 $A_{1} \times \cdots \times A_{m}$ 恰为 $X$ 的所有 $m$-串的集合，记作 $X^{m}$。类似地，笛卡儿积 $A_{1} \times A_{2} \times \cdots$ 恰为 $X$ 的所有 $\omega$-串的集合，记作 $X^{\omega}$。

后面我们还将对任意加标集族定义笛卡儿积。

\begin{example}\label{example:ch01-欧氏空间}
设 $\mathbb{R}$ 为实数集，则 $\mathbb{R}^{m}$ 表示实数的所有 $m$-串的集合，常称为 $m$-维欧氏空间（Euclidean $m$-space）（虽然欧几里得从未研究过它）。类似地，$\mathbb{R}^{\omega}$ 有时称为“无穷维欧氏空间”，它是所有实数 $\omega$-串 $(x_{1}, x_{2}, \cdots)$ 的集合，即所有函数 $\boldsymbol{x}: \mathbb{Z}_{+} \longrightarrow \mathbb{R}$ 的集合。
\end{example}

\subsection*{习题}

\begin{exercise}\label{exercise:ch01-笛卡儿积-1}
证明：存在 $A \times B$ 与 $B \times A$ 之间的一个一一对应。
\end{exercise}

\begin{exercise}\label{exercise:ch01-笛卡儿积-2}
(a) 证明：对于 $n > 1$，存在
\begin{align*}
A_{1} \times \cdots \times A_{n} \quad \text{与} \quad (A_{1} \times \cdots \times A_{n-1}) \times A_{n}
\end{align*}
之间的一个一一对应。

(b) 给定一个加标集族 $\{A_{1}, A_{2}, \cdots\}$，对于每一正整数 $i$，记 $B_{i} = A_{2i-1} \times A_{2i}$。证明存在 $A_{1} \times A_{2} \times \cdots$ 与 $B_{1} \times B_{2} \times \cdots$ 之间的一个一一对应。
\end{exercise}

\begin{exercise}\label{exercise:ch01-笛卡儿积-3}
设 $A = A_{1} \times A_{2} \times \cdots$，$B = B_{1} \times B_{2} \times \cdots$。
\begin{enumerate}[(a)]
\item 证明：若对于每一个 $i$，有 $B_{i} \subset A_{i}$，则 $B \subset A$。（严格地讲，如果我们给出的是从 $\mathbb{Z}_{+}$ 到所有 $B_{i}$ 之并的一个函数，那么在把它作为从 $\mathbb{Z}_{+}$ 到所有 $A_{i}$ 之并的一个函数处理时，必须先改变其值域。当研究笛卡儿积时，这一点可以忽略。）
\item 证明：若 $B$ 为非空集合，则 (a) 的逆命题成立.
\item 证明：若 $A$ 为非空集合，则 $A_{i}$ 为非空集合. 其逆命题成立吗？（在第 19 节的练习中，我们将再度提到这个问题.）
\item $A \cup B$ 与 $A_{i} \cup B_{i}$ 的笛卡儿积之间的关系？$A \cap B$ 与 $A_{i} \cap B_{i}$ 的笛卡儿积之间的关系？
\end{enumerate}
\end{exercise}

\begin{exercise}\label{exercise:ch01-笛卡儿积-4}
设 $m$，$n \in \mathbb{Z}_{+}$ 并且 $X \neq \varnothing$。
\begin{enumerate}[(a)]
\item 若 $m \leqslant n$，给出一个单射 $f: X^{m} \longrightarrow X^{n}$.
\item 给出一个一一对应 $g: X^{m} \times X^{n} \longrightarrow X^{m+n}$.
\item 给出一个单射 $h: X^{n} \longrightarrow X^{\omega}$.
\item 给出一个一一对应 $k: X^{n} \times X^{\omega} \longrightarrow X^{\omega}$.
\item 给出一个一一对应 $l: X^{\omega} \times X^{\omega} \longrightarrow X^{\omega}$.
\item 若 $A \subset B$，给出一个单射 $m: (A^{\omega})^{n} \longrightarrow B^{\omega}$.
\end{enumerate}
\end{exercise}

\begin{exercise}\label{exercise:ch01-笛卡儿积-5}
$\mathbb{R}^{\omega}$ 的下列子集中哪些能够表示成 $\mathbb{R}$ 的子集的笛卡儿积？
\begin{enumerate}[(a)]
\item $\{\boldsymbol{x} \mid \text{对于所有 } i, x_{i} \text{ 为整数}\}$.
\item $\{\boldsymbol{x} \mid \text{对于所有 } i, x_{i} \geqslant i\}$.
\item $\{\boldsymbol{x} \mid \text{对于所有 } i \geqslant 100, x_{i} \text{ 为整数}\}$.
\item $\{\boldsymbol{x} \mid x_{2} = x_{3}\}$.
\end{enumerate}
\end{exercise}

\end{document}